\section{Introduction}

\subsection{The Static-Content Problem}

Educational content has historically been authored once and revised
infrequently. A textbook chapter, a worked example, or an interactive diagram
is designed using the author's best judgment about what will help students
learn, and then deployed largely unchanged for years. This stands in stark
contrast to how learning actually works: the effectiveness of any given
explanation varies enormously across learners, contexts, and prior knowledge
states.

The emergence of \textit{intelligent textbooks}---interactive, AI-assisted
learning materials that incorporate concept dependency graphs, learning-objective
taxonomies, and embedded simulations---creates an opportunity to break this
static mold. Because these textbooks are instrumented (they can record how
students interact with content) and modular (individual components can be
regenerated independently), they are uniquely positioned to support
\textit{continuous, evidence-driven improvement} of their own pedagogical
artifacts.

\subsection{Learning from Self-Improving AI}

A parallel revolution has been unfolding in artificial intelligence research.
A class of systems has emerged that can improve their own performance through
iterative self-modification, validated against empirical benchmarks rather than
human design alone. The most relevant of these is the
\textbf{Darwin-G\"{o}del-Machine (DGM)}, introduced in 2025 by researchers at
Sakana AI, the University of British Columbia, and the Vector
Institute~\cite{zhang2025dgm,sakana2025dgm}.

The central insight of the DGM is a pragmatic relaxation. The original
G\"{o}del Machine, proposed theoretically by J\"{u}rgen
Schmidhuber~\cite{schmidhuber2007goedel}, required that any self-modification
be \textit{provably} beneficial before adoption---a requirement that is
intractable for complex systems. The DGM instead requires only
\textit{empirical evidence} that a proposed modification improves performance
on real benchmarks. This single shift transforms self-improvement from a
theoretical impossibility into an engineering practice.

This paper argues that the same shift can be applied to educational content.
We do not need to \textit{prove} that a given MicroSim or explanation is
pedagogically optimal. We need only gather empirical evidence---from learning
analytics---that one variant outperforms another, and then let an archive of
variants evolve over time. The remainder of this paper develops this analogy
into a concrete, implementable framework.

\subsection{Contributions}

This paper makes the following contributions:

\begin{enumerate}
    \item A conceptual mapping between DGM mechanisms and intelligent textbook
          architecture.
    \item A survey of the current state of the art in self-improving AI systems
          relevant to education.
    \item A worked architecture using MkDocs, p5.js MicroSims, and an
          xAPI/LRS analytics pipeline.
    \item A proposed JSON schema standard for storing MicroSim performance
          data.
    \item A discussion of safety constraints and human oversight requirements
          for evolving educational content.
\end{enumerate}
